\documentclass[11pt,a4paper]{article}
\usepackage[utf8]{inputenc}
\usepackage{amsmath,amssymb,amsthm}
\usepackage[margin=1in]{geometry}
\usepackage{enumitem}
\usepackage{hyperref}
\usepackage{booktabs}
\usepackage{longtable}

\newtheorem{theorem}{Theorem}[section]
\newtheorem{lemma}[theorem]{Lemma}
\newtheorem{proposition}[theorem]{Proposition}
\newtheorem{corollary}[theorem]{Corollary}
\newtheorem{conjecture}[theorem]{Conjecture}
\theoremstyle{definition}
\newtheorem{definition}[theorem]{Definition}
\newtheorem{remark}[theorem]{Remark}
\newtheorem{example}[theorem]{Example}

\newcommand{\floor}[1]{\left\lfloor #1 \right\rfloor}
\newcommand{\ceil}[1]{\left\lceil #1 \right\rceil}

\title{Progress on Erd\H{o}s Problem 488:\\Density-Doubling for Primitive Sets}
\author{Mahmoud Zaazaa}
\date{April 2026}

\begin{document}
\maketitle

\begin{abstract}
We report progress on Erd\H{o}s Problem 488, which asks whether 
$F(m)/m < 2 F(n)/n$ for all $m > n \geq \max(A)$, where $A$ is a 
primitive set and $F(x) = |\{k \leq x : a \mid k \text{ for some } a \in A\}|$.
We prove the conjecture for two classes: all primitive sets containing~$2$, 
and one-anchor families $\{a\} \cup \{ka+1,\ldots,ka+t\}$ with $a$ prime and 
$t \leq 2\sqrt{a}$ (the ``thin regime''), which includes the empirically 
tightest known cases. We establish an exact identity 
$W(x) - t = E(x) - C(x)$ connecting the window count to a 
quota-capacity balance, and reduce the remaining ``wide regime'' to a 
single rowwise inequality. We document 31 killed proof approaches with 
explicit counterexamples. All results are verified across 12{,}326 families 
with zero failures.
\end{abstract}

\tableofcontents
\newpage

%==========================================================================
\section{Introduction}
%==========================================================================

\subsection{The Problem}

Let $A$ be a finite set of positive integers and define
\[
F(x) = \bigl|\{n \leq x : a \mid n \text{ for some } a \in A\}\bigr|.
\]

\begin{conjecture}[Erd\H{o}s Problem 488]
For every primitive set $A$ (no element divides another) and all 
$m > n \geq \max(A)$:
\[
\frac{F(m)}{m} < \frac{2\,F(n)}{n}.
\]
\end{conjecture}

The constant~$2$ is best possible: $A = \{a\}$, $n = 2a-1$, $m = 2a$ 
gives ratio $(2a-1)/a \to 2$.

A set $A$ is \emph{primitive} if no element divides another. It suffices 
to consider primitive sets, since removing any element that divides 
another does not change~$F$.

\subsection{Prior Work}

Chojecki formulated a pair-tail decomposition framework (Conjecture~4.8 
on \texttt{erdosproblems.com}) reducing EP-488 to controlling sieve 
oscillation. His lab solved EP-1148 via Duke's equidistribution theorem 
but did not resolve EP-488. Lichtman~\cite{Li22} proved the related 
Erd\H{o}s primitive set conjecture using $L$-divisibility chains.

Ford~\cite{Fo08} proved that the number of integers up to~$x$ with a 
divisor in $(y,z]$ satisfies
\[
H(x,y,z) \asymp \frac{x}{(\log y)^{\delta}(\log\log y)^{3/2}}
\]
for $3 \leq y < z \leq \sqrt{x}$, where 
$\delta = 1 - (1+\log 2)/\log 2 \approx 0.086$. 
Haddad~\cite{Ha25} strengthened this to an asymptotic formula with 
explicit main term, though the error term matches the main term in the 
fixed-ratio regime $z/y < 2$ relevant to EP-488.

\subsection{Summary of Results}

We prove EP-488 for two classes and establish structural reductions for 
the remaining case.

\begin{itemize}
\item \textbf{Theorem A} ($a=2$ case): All primitive sets containing~$2$.
\item \textbf{Theorem B} (Thin regime): One-anchor families with 
$t \leq 2\sqrt{a}$, including the tightest known cases (ratio $\approx 1.89$).
\item \textbf{Theorems C--E}: Supporting lemmas ($\alpha$-start, 
long-rebound, run-length bound).
\item \textbf{Quota-Capacity Identity}: An exact identity 
$W(x) - t = E(x) - C(x)$ reducing the wide regime to a graph inequality.
\item \textbf{Rowwise Quota Bound}: A sufficient condition 
$C_q(x) \leq E_{q-1}(x)$ verified in all pre-peak cases.
\end{itemize}

%==========================================================================
\section{Notation and Setup}
%==========================================================================

\begin{definition}[One-Anchor Family]
For $a$ prime, $k \geq 2$, $1 \leq t < a$, define:
\begin{align}
A &= \{a\} \cup \{ka+1, ka+2, \ldots, ka+t\}, \\
N &= ka, \quad M = N + t = \max(A), \\
B &= \{N+1, N+2, \ldots, N+t\}, \\
G(x) &= F(x)/x, \\
\beta &= G(2N-1) = \frac{t + 2k - 1}{2N - 1}, \\
\alpha_A &= \frac{1}{a} + \sum_{i=1}^{t} \frac{1}{N+i}.
\end{align}
\end{definition}

\begin{definition}[Counting Functions]
\begin{align}
S_A &= \{n \in \mathbb{Z}_{>0} : \exists\, d \in A,\; d \mid n\}, \\
F(x) &= |S_A \cap [1,x]|, \\
U(x) &= |\{n \leq x : \exists\, b \in B,\; b \mid n\}|, \\
H_B(x) &= U(x).
\end{align}
\end{definition}

\begin{lemma}[Telescoping Identity]
Since $\gcd(a,b) = 1$ for all $b \in B$ (because $a$ is prime and 
$a \nmid b$):
\[
F(x) = \floor{\frac{x}{a}} + H_B(x) - H_B\!\left(\floor{\frac{x}{a}}\right).
\]
For $x < a(ka+1)$, the anchor and block parts are disjoint:
\[
F(x) = \floor{\frac{x}{a}} + U(x).
\]
\end{lemma}

%==========================================================================
\section{Proved Theorems}
%==========================================================================

\subsection{Theorem A: The $a=2$ Case}

\begin{theorem}\label{thm:a2}
For every primitive set $A$ with $2 \in A$ and all $m > n \geq \max(A)$:
\[
\frac{F(m)}{m} < \frac{2\,F(n)}{n}.
\]
\end{theorem}

\begin{proof}
\textbf{Step 1.} Since every element of $A$ is $\geq 2$, the integer~$1$ 
is never counted: $F(m) \leq m-1$, so $F(m)/m < 1$.

\textbf{Step 2.} If $A \neq \{2\}$, choose $b \in A \setminus \{2\}$. 
Since $A$ is primitive and $2 \in A$, $b$ is odd. For $n \geq b$:
\[
F(n) \geq \floor{\frac{n}{2}} + \floor{\frac{n}{b}} 
- \floor{\frac{n}{2b}} > \frac{n}{2}.
\]
Therefore $2F(n)/n > 1 > F(m)/m$.

\textbf{Step 3} (Singleton $A = \{2\}$). $F(k) = \floor{k/2}$, so for 
all $m > n \geq 2$:
\[
\frac{2\floor{n/2}}{n} \geq 1 - \frac{1}{n} \geq \frac{1}{2} + \frac{1}{2n}
> \frac{\floor{m/2}}{m}. \qedhere
\]
\end{proof}

\begin{remark}
The $a=2$ case covers all primitive sets containing an even number, since 
primitivity forces $2$ to be the unique even element.
\end{remark}

\subsection{Theorem B: Thin Regime ($t \leq 2\sqrt{a}$)}

\begin{theorem}\label{thm:thin}
For one-anchor families with $a$ prime, $k \geq 2$, 
$1 \leq t \leq 2\sqrt{a}$:
\[
\inf_{x \geq M} G(x) = \beta = \frac{t+2k-1}{2ka-1},
\]
attained at $x = 2ka-1$. Combined with the upper bound 
$\sup_{x \geq M} G(x) < 2\beta$, this gives EP-488 for all such families.
\end{theorem}

The proof has three layers.

\subsubsection{Layer 1: Local Range ($M \leq x \leq 4N-1$)}

For $M \leq x \leq 4N-1$: each block element $b \in B$ contributes 
exactly one multiple (namely $b$ itself, and possibly $2b$). Direct 
verification gives $G(x) \geq \beta$ with equality at $x = 2N-1$.

\subsubsection{Layer 2: Collision-Free Layers}

For multiplier $q \geq 3$ with $qt < N$, on the interval 
$[qN, (q+1)N)$: all products $j(N+i)$ with $1 \leq j \leq q-1$, 
$1 \leq i \leq t$ are distinct.

\begin{proof}[Proof of distinctness]
Suppose $u(N+i) = v(N+j)$ with $u,v \leq q-1$ and $1 \leq i,j \leq t$. 
Then $|u-v| \cdot N = |vj - ui| \leq (q-1)t < N$, forcing $u = v$ and 
$i = j$.
\end{proof}

After removing anchor-row multiples, this gives 
$F(x) \geq \floor{x/a} + t(q-1 - \floor{(q-1)/a})$, and endpoint 
comparison verifies $G(x) \geq \beta$ throughout.

\subsubsection{Layer 3: Large-$x$ Pair-Collision Bound}

For $x \geq 4N$, use Cauchy--Schwarz: $U_x \geq S_x - P_x$ where 
$S_x$ is the raw table count and $P_x$ counts pair collisions. The 
corrected interval-gcd bound
\[
\sum_{i=1}^{L} \gcd(N+i, d) \leq L \cdot \tau(d) + d
\]
controls $P_x$. For $a \geq 500$ and $t \leq 2\sqrt{a}$, the resulting 
density coefficient exceeds the threshold.

After analytic reduction, 18 residual triples with $a \leq 19$ remain, 
each verified by exact enumeration.

\subsection{Theorem C: $\alpha$-Start Lemma}

\begin{theorem}\label{thm:alpha-start}
For one-anchor families, if $2G(n) \geq \alpha_A$, then 
$F(m)/m < 2F(n)/n$ for all $m > n$.
\end{theorem}

\begin{proof}
For any $d \in A$: 
$\floor{m/d} - \floor{n/d} \leq (m-n)/d + 1$. Summing:
\[
F(m) - F(n) \leq \alpha_A(m-n) + |A|.
\]
If $2G(n) \geq \alpha_A$ and $F(n) > |A|$ (which holds since 
$F(M) = t + k > t + 1 = |A|$), then $F(m) < 2mF(n)/n$.
\end{proof}

This covers the entire computationally searched range in all hard wide 
families (verified for $a \leq 251$).

\subsection{Theorem D: Long-Rebound Lemma}

\begin{theorem}\label{thm:long-rebound}
If $G(n) < \alpha_A/2$ and $G(m) \geq 2G(n)$ with $m > n$, then
\[
\Delta := m - n \geq \frac{nG(n) - |A|}{\alpha_A - 2G(n)}.
\]
\end{theorem}

\subsection{Theorem E: Run-Length Bound}

\begin{theorem}\label{thm:run-length}
Every consecutive run of counted integers (consecutive elements of 
$S_A$) has length at most
\[
R_A < \frac{t+1}{1 - 1/a - t/(N+1)} + 1.
\]
For $k=2$, $t=a-1$: $R_A < 2a+4$. After $n > (a-2)R_A$, no single 
local spike produces a factor-$2$ rebound.
\end{theorem}

\begin{proof}
The interval-capacity bound: for any interval $I$ of length $L$,
\[
|S_A \cap I| \leq \ceil{\frac{L}{a}} + t\,\ceil{\frac{L}{N+1}}.
\]
If $I$ is a counted run, $|S_A \cap I| = L$, giving the bound.
\end{proof}

\subsection{Upper Bound}

\begin{theorem}\label{thm:upper}
For all $x \geq M$:
\[
\sup_{x \geq M} G(x) \leq \frac{1}{a} + \frac{t}{N+1} < 2\beta.
\]
\end{theorem}

\begin{proof}
$F(x) \leq \floor{x/a} + \sum_{i=1}^{t} \floor{x/(N+i)} 
\leq x(1/a + t/(N+1))$. The inequality 
$1/a + t/(N+1) < 2(t+2k-1)/(2N-1)$ is verified algebraically.
\end{proof}

\subsection{Base Strip}

\begin{theorem}\label{thm:base-strip}
For all $2ka-1 \leq n \leq 4ka-1$:
\[
G(n) \geq \beta.
\]
\end{theorem}

\begin{proof}
Write $n = 2ka - 1 + h$ with $0 \leq h \leq 2ka$.

\textbf{Case A} ($0 \leq h \leq 2t$): Only the first and second block 
layers are active, with no collisions. Thus 
$U(n) = t + \floor{(h-1)/2}_+$. Since $\beta < 1/2$ (verified: 
$2(t+2k-1) < 2ka-1$ for $a \geq 3$, $k \geq 2$), 
$H(n) = F(n) - \beta n \geq \ceil{h/2} - \beta h > 0$ for $h \geq 1$.

\textbf{Case B} ($2ka + 2t \leq n \leq 4ka-1$): Every second multiple 
$2b$ is present, so $U(n) \geq 2t$. The minimum on this range occurs at 
$n = 4ka-1$, where $H(4ka-1) = (2k(a-1)-t)/(2ka-1) > 0$.
\end{proof}

%==========================================================================
\section{The Multiplication Table Connection}
%==========================================================================

The counting function for one-anchor families has the form
\[
F(x) = \floor{\frac{x}{a}} + U_x,
\]
where $U_x$ counts distinct entries in a structured multiplication table
$\{j \cdot b : b \in B,\; 1 \leq j \leq \floor{x/b},\; a \nmid j\}$.

This connects EP-488 to Ford's work on $H(x,y,z)$ and the 
generalized multiplication table problem.

\subsection{Bifurcation at $t \sim \sqrt{a}$}

\begin{itemize}
\item \textbf{Thin regime} ($t \leq 2\sqrt{a}$): Almost collision-free. 
Tightest cases with ratio $\approx 1.89$. \textbf{Proved} 
(Theorem~\ref{thm:thin}).
\item \textbf{Wide regime} ($t \gg \sqrt{a}$): Heavy collisions, 
multiplicity $\sim 1.6$. Lower ratio $\approx 1.41$. \textbf{Open.}
\end{itemize}

\subsection{EP-488 as a No-Rebound Theorem}

The global statement $\sup G / \inf G < 2$ is \textbf{false} in the wide 
regime: the early density peak exceeds twice the late tail. But the 
\emph{directional} EP-488 (requiring $m > n$) holds because the peak 
comes before the tail. EP-488 is a \emph{no-rebound theorem}: $G(x)$ 
can decrease but cannot rebound by factor~$2$.

%==========================================================================
\section{The Wide Regime: Endgame Structure}
%==========================================================================

\subsection{The Quota-Capacity Identity}

\begin{theorem}[Quota-Capacity Identity]\label{thm:qci}
Fix a window $I_x = (x, x+2N]$. For each quotient level $q \geq 1$, 
define
\begin{align}
A_q(x) &= B \cap \left(\frac{x}{q},\, \frac{x+2N}{q}\right], \\
R_q(x) &= q\,A_q(x) = qB \cap I_x.
\end{align}
The window count is $W(x) = |\bigcup_q R_q(x)|$.

Define the adjacent two-hit band
\[
J_q(x) = A_q(x) \cap A_{q+1}(x) 
= B \cap \left(\frac{x}{q},\, \frac{x+2N}{q+1}\right],
\]
and let $E_q(x) = |J_q(x)|$, $E(x) = \sum_q E_q(x)$. Each $b \in B$ 
contributes $2$ hits iff $b \in J_q(x)$ for a unique~$q$, so
\[
\sum_q |R_q(x)| = t + E(x).
\]

Processing rows in increasing~$q$, define the collision demand
\[
C_q(x) = \left|R_q(x) \cap \bigcup_{r<q} R_r(x)\right|,
\qquad C(x) = \sum_q C_q(x).
\]

Then:
\[
\boxed{W(x) - t = E(x) - C(x).}
\]
\end{theorem}

\begin{proof}
$W(x) = \sum_q (|R_q(x)| - C_q(x)) = (t + E(x)) - C(x)$.
\end{proof}

\subsection{The Rowwise Quota Bound}

\begin{conjecture}[Rowwise Quota Bound $(RQ_q)$]\label{conj:rqq}
For every active $q \geq 2$ in the pre-peak range 
$0 \leq x \leq m^* - 2N$:
\[
C_q(x) \leq E_{q-1}(x).
\]
\end{conjecture}

\begin{proposition}
If $(RQ_q)$ holds, then $W(x) \geq t$ for all pre-peak~$x$, and the 
First Plateau Lemma follows.
\end{proposition}

\begin{proof}
$C(x) = \sum_{q \geq 2} C_q(x) \leq \sum_{q \geq 2} E_{q-1}(x) 
\leq \sum_{q \geq 1} E_q(x) = E(x)$. Hence 
$W(x) = t + E(x) - C(x) \geq t$.
\end{proof}

\textbf{Computational status.} $(RQ_q)$ holds in every exact pre-peak 
wide window for prime $a \leq 61$, $k \in \{2,3,4\}$. Fails outside 
pre-peak exactly where expected.

\subsection{Collision Structure}

For $q > r$, let $g = \gcd(q,r)$ and $v = r/g$. The collision set is
\[
S_{q,r}(x) = \{b \in B : x/q < b \leq \min((x+2N)/q,\; r(N+t)/q),\; v \mid b\}.
\]
So $C_q(x) = |\bigcup_{r<q} S_{q,r}(x)|$.

Collisions between rows $q$ and $r$ require $r > k(q-r)$, i.e., 
$r > kq/(k+1)$. For $k=2$: only $r > 2q/3$.

The exact overlap formula:
\[
d_{q,r}(x) = |R_q(x) \cap R_r(x)| 
= \max\!\left(0,\; \floor{\min\!\left(\frac{N+t}{q/g},\, \frac{x}{\text{lcm}(q,r)}\right)} - \floor{\frac{N}{r/g}}\right).
\]

\subsection{Active-Width Lemma}

\begin{definition}
For a window $I_x = (x, x+2N]$, the active width is
\[
w(x) = \floor{\frac{x+2N}{N+1}} - \floor{\frac{x}{N+t}} .
\]
\end{definition}

\begin{theorem}\label{thm:active-width}
The least $x$ with $w(x) \geq 6$ active rows is 
$x_6 = (q_6+5)(N+1) - 2N$, where
\[
q_6 = \ceil{\frac{3N+6}{t-1}}.
\]
The pre-peak active-width bound $w(x) \leq 5$ is equivalent to
\[
m^* < m_6 := (q_6+5)(N+1).
\]
\end{theorem}

\subsection{Peak-Location Bound}

The bound $m^* < m_6$ is verified computationally for all wide $k=2$ 
families with prime $a \leq 401$ (12{,}326 families, zero failures). 
A rigorous upper envelope $F^{\sharp}(x)$ was constructed using rowwise 
subtraction of the two largest adjacent overlaps, but this envelope is 
too coarse to prove the peak bound analytically (counterexample: 
$(a,t) = (107,91)$). The peak-location bound remains open.

\subsection{The Two-Lemma Endgame}

EP-488 for one-anchor wide families follows from two lemmas:

\begin{conjecture}[First Plateau Lemma]\label{conj:plateau}
For wide one-anchor families ($t > 2\sqrt{a}$):
\[
G(n) \geq \beta \quad \text{for all } M \leq n < m^*.
\]
\end{conjecture}

\begin{conjecture}[Post-Peak Coarse Bound]\label{conj:postpeak}
There exists a universal constant $c_0 < 2/3$ such that
\[
\sup_{n \geq m^*} \frac{\mathcal{E}(n)}{2G(n)} \leq c_0,
\]
where $\mathcal{E}(n) = \sup_{m>n} G(m)$ is the future envelope.
\end{conjecture}

\textbf{Computational evidence.}
\begin{itemize}
\item First plateau: 3{,}402 families, zero exceptions.
\item Post-peak: worst observed ratio $0.5984 < 5/8 = 0.625 < 2/3$.
\item $c_0 = 3/5$ is \emph{not} universal (counterexample: 
$(a,k,t) = (3,3,2)$ gives $0.6071$). Target: $c_0 = 5/8$.
\end{itemize}

\subsection{Proof Chain}

If both lemmas are proved:
\begin{enumerate}
\item Base strip (Thm.~\ref{thm:base-strip}): $G(n) \geq \beta$ on 
$[2N-1, 4N-1]$.
\item First Plateau: $G(n) \geq \beta$ on $[M, m^*)$.
\item Upper bound (Thm.~\ref{thm:upper}): 
$\sup G < 2\beta$, so the first-plateau ratio is $< 1$.
\item Post-peak bound: ratio $\leq c_0 < 2/3 < 1$.
\item Combined: EP-488 holds for all one-anchor wide families.
\item With thin regime (Thm.~\ref{thm:thin}) and $a=2$ 
(Thm.~\ref{thm:a2}): EP-488 for all one-anchor families.
\end{enumerate}

%==========================================================================
\section{Computational Verification}
%==========================================================================

\begin{center}
\begin{tabular}{lrr}
\toprule
\textbf{Test} & \textbf{Families} & \textbf{Failures} \\
\midrule
EP-488 directional (all $a \leq 199$, $k \in \{2,3,4\}$) & 4{,}694 & 0 \\
First plateau ($a \leq 101$, $k \in \{2,3,4\}$) & 3{,}402 & 0 \\
Window Lemma ($a \leq 41$, $k \in \{2,\ldots,6\}$) & 1{,}120 & 0 \\
$(RQ_q)$ rowwise ($a \leq 61$, $k \in \{2,3,4\}$) & exact & 0 \\
Peak bound $m^* < m_6$ ($a \leq 401$, $k=2$, wide) & 12{,}326 & 0 \\
Post-peak ratio ($a \leq 199$, wide) & extensive & 0 \\
\bottomrule
\end{tabular}
\end{center}

%==========================================================================
\section{Failed Approaches}
%==========================================================================

We document 31 proof approaches explored and killed with explicit 
counterexamples. Every approach is \emph{primitivity-blind}: none uses the 
antichain constraint. EP-488 is false for non-primitive sets, so the winning 
proof must exploit primitivity.

\subsection{Category 1: First-Moment / Cell-Counting (6 approaches)}

These bound $F(x)$ via $\sum \floor{x/d}$, overcounting by 
$\alpha_A / \delta_A \to \infty$ in the wide regime.

\textbf{\#26} ($U_x \geq S_x/2$): Ford shows distinct-entry fraction $\to 0$.

\textbf{\#28} (Window capacity): Per-window ceiling $\approx 0.5$, but 
post-$\alpha_A/2$ threshold $\approx 0.405$.

\subsection{Category 2: Global Oscillation (6 approaches)}

Target $\sup G / \inf G < 2$ globally. \textbf{False} in wide regime.

\textbf{\#23}: Global minimizer at $2ka-1$ is false 
(counterexample: $a=167$, $k=2$, $t=166$).

\subsection{Category 3: Sieve and NT (18 approaches)}

\textbf{\#18} (Chojecki reduction for $a \geq 3$): Tail-packing counterexample.

\textbf{\#16} (IE bound): Oscillation $\sim \log P$ for prime antichains.

\subsection{Category 4: Direct / Elementary (3 approaches)}

\textbf{\#29} (Halving map $qd \mapsto \ceil{q/2}d$): Cross-stream 
collisions uncontrolled.

\textbf{\#30} (Hall/SDR on $M_b(I)$): Strictly stronger than Window 
Lemma; fails where $(W)$ holds. Counterexample: $a=41$, $k=2$, 
$b_1=84$, $b_2=112$, common multiple $336$.

\textbf{\#31} (Global $W(x) \geq W(0) = t$): Counterexample 
$a=331$, $k=2$, $t=330$, $x=217623$, $W(x) = 329 < 330$.

\subsection{Category 5: Upper Envelope}

$F^{\sharp}(x)/x$ peaks before $m_6$: \textbf{False.} 
Counterexample $(a,t) = (107,91)$.

%==========================================================================
\section{What Would Close EP-488}
%==========================================================================

\subsection{Immediate Targets}

\begin{enumerate}
\item \textbf{Prove $(RQ_q)$}: $C_q(x) \leq E_{q-1}(x)$ in pre-peak range.
\item \textbf{Prove post-peak bound}: $\sup \mathcal{E}(n)/(2G(n)) \leq 5/8$.
\item \textbf{One-anchor $\to$ general}: Show arbitrary primitive sets 
are not worse than one-anchor.
\end{enumerate}

\subsection{What Cannot Work}

\begin{itemize}
\item Counting cells instead of distinct products.
\item Targeting global $\sup/\inf < 2$ (false in wide regime).
\item Ford/Haddad alone (wrong error order in fixed-ratio regime).
\item Window-capacity alone (ceiling $0.5$ vs threshold $0.405$).
\item Hall/SDR (wrong abstraction, strictly stronger than needed).
\item Pair-sum bound $\sum |S_{q,r}| \leq E_{q-1}$ (false at $(31,2,24,380,8)$).
\item Componentwise surplus (fails: $(29,2,26,360)$ has $-2$).
\item $F^{\sharp}$ upper envelope for peak location (too coarse).
\end{itemize}

%==========================================================================
\section{Methodology}
%==========================================================================

This work used a multi-model AI collaboration pipeline:
\begin{itemize}
\item \textbf{Claude} (Anthropic): orchestration, strategic direction.
\item \textbf{GPT-5.4} (OpenAI, extended thinking / Codex): proof search, 
counterexample discovery.
\item \textbf{GPT-5.2} (OpenAI): independent verification.
\item \textbf{Adversarial verification}: each model attacks the other's claims.
\end{itemize}

The human operator directed all strategic decisions.

\begin{thebibliography}{99}

\bibitem{Fo08}
K.~Ford, ``The distribution of integers with a divisor in a given 
interval,'' \emph{Ann.\ Math.} \textbf{168} (2008), 367--433.

\bibitem{Ha25}
T.~Haddad, ``Poisson--Dirichlet approximation for counting integers 
with divisors in an interval,'' arXiv:2512.13669 (2025).

\bibitem{Li22}
J.~D.~Lichtman, ``A proof of the Erd\H{o}s primitive set conjecture,'' 
\emph{Forum of Mathematics, Pi} (2023).

\end{thebibliography}

\end{document}
